\subsection{Scope of the Finite-Group Window Comparisons}
  \label{subsec:SU2-thread-status}

  The embeddings $Q_8\subset2I\subset SU(2)$ relate finite-group examples to a continuous group.
  They do not supply a convergence theorem or a physical attractor selected by the rate bound.

  Under the single-mode and linearised-dispersion hypotheses of
  Section~\ref{subsec:admissibility-bounded-flux}, ratios of maximal effective amplitudes obey
  \begin{equation}
    \frac{A_\rho^{\max}}{A_\sigma^{\max}}
    =\sqrt{\frac{\lambda_\sigma}{\lambda_\rho}},\qquad
    \lambda_\rho,\lambda_\sigma>0.
  \end{equation}
  For $Q_8$ with $S=\{\pm i,\pm j,\pm k\}$, the two-dimensional sector relative to a non-trivial
  one-dimensional sector gives $\sqrt{4/3}$.
  For $2I$ with $S=10a\cup10b$, $|S|=24$, the spin-$1/2$ window relative to the spin-$5/2$ window
  gives $\sqrt{14/9}$~\cite{Beau2026a1}.
  These denominators are different objects; $2I$ has no non-trivial abelian sector.
  Neither contrast selects the occupied sector or implies saturation.

  LPS graphs do not form an intermediate subgroup in these embeddings.
  No corresponding sector ratio or convergence of it to unity is specified here for LPS or $SU(2)$.
  Such a continuum statement requires a common comparison and a convergence construction.
  The capacity criterion of Section~\ref{subsec:binary-polyhedral-maximality} is a separate proposal;
  it does not follow from these ratios or establish an attractor from [A-cap].
